%% kapa: exam=1 round=26 year=2015 period=1 form=A subject=economics subject_ko=경제원론 seq=5 number=45
%% source: Q-Net 감정평가사 기출문제 (cst00309) artlSeq=5138001 file=extracted/1008/5138001/제26회_감정평가사 1차_기출/제26회_감정평가사 1차_기출/1교시_감평_A형.pdf
%% conversion: kordoc 4.24.51 -> tools/gen1.py (auto); review=proofread-sonnet; flags=none
\begin{question}{45}
甲의 효용함수는 $U(x,y)=xy$이고, $X$재와 $Y$재의 가격이 각각 1과 2이며, 甲의 소득은 100이다. 예산제약 하에서 甲의 효용을 극대화시키는 $X$재와 $Y$재의 소비량은? (단, 甲은 $X$재와 $Y$재만 소비하고 $x$는 $X$재의 소비량, $y$는 $Y$재의 소비량이다.)
\begin{choices}
  \item $x=20, y=40$
  \item $x=30, y=35$
  \item $x=40, y=30$
  \item $x=50, y=25$
  \item $x=60, y=20$
\end{choices}
\end{question}
\answer{4}
